% !TeX root = ../papers/nc-p3.tex
\section{種数一の被約連結成分と五角形の非存在}
\label{sec:genus-one-polygon-exclusion}
\begin{theorem}\label{thm:reduced-genus-one-components}
点ファイバーを持つ標準有限Frobenius窓の共通点グラフで、純な被約曲線が開閉な連結成分 $C$ をなすとする。
$p_a(C)=0$ は不可能である。$p_a(C)=1$ なら滑らかな楕円曲線か、滑らかな有理成分を
二つずつの節点で輪につないだ多角形である。後者では四移動が $C$ を保つときは固定点のない回転を誘導する。
$C$ が点グラフの唯一の曲線成分をなす、という仮定はここでいう連結成分の分類には不要である。
\end{theorem}
\begin{proof}
正規化の成分と特異点を頂点とし、各特異点への各枝を辺とする連結二部多重グラフを取る。
正規化の完全列から
\[
 p_a(C)=\sum g(\widetilde C_j)+b_1(\text{グラフ})+
            \sum_p(\delta_p-r_p+1)
\]
である。$r_p$ は枝数で、最後の和の各項は非負である。
実際、完備局所環は正規化の各枝の定数項だけを一致させた環 $A'$ に入り、
この項は $A'$ までの余長となる。余長零の特異点は $r$ 本の独立な座標軸の和であり、
埋込み次元上限から $r=2,3$ である。以下、$C$ が他の連結成分に移れば一致点を持たないことにも注意する。

木の末端にある有理正規化成分がただ一つの上記の座標軸特異点で会う場合を考える。
この枝では $B=0$ で余接捩れは $r-1$ である。局所式 $x_i x_j=0$ を一枝上で微分すればよい。
定理\ref{thm:fixed-differential}の恒等の場合は $r-1=4$ が必要であり、非恒等の場合に一致因子の長さは
$r-3$ となる。$r=2$ では負、$r=3$ では命題\ref{prop:automatic-spatial-tangent-mark}により
その点が四移動で固定されるため長さ零に反する。滑らかな開閉有理成分も同じ式で除ける。
従って種数零の木は消える。種数一で正規化に楕円成分があれば、木に他の成分があると
この楕円成分以外の末端を持つので、滑らかな楕円だけとなる。
$b_1=1$ の場合は全て座標軸特異点で、同じ末端排除により全頂点の価数は2以上、Euler式から全て2である。
従って節点の輪となる。これが四移動で保たれるとき各有理成分を点ごとに固定することはできず、
一致因子の長さは各成分で零である。向きの反転は節点か成分を保ち、後者でも射影直線上に固定点を持つ。
同様に一成分の自己節点の輪も、唯一の節点が固定されるため保たれる場合は生じない。

残る場合は木にただ一つ、$\delta_p-r_p+1=1$ の特異点 $p$ がある。
$p$ から遠い方へ辿れば、$p$ に直接付く成分の外には先に除いた末端しか作れない。
従って全成分は $p$ の一枝ずつに対応する有理成分で、他の特異点はない。
この場合 $C$ が保たれるなら $p$ は固定される。保たれない可能性については次数上限も使う。
$C$ とその像が相異なれば両方の次数は同じで合計7以下なので $\deg C\le3$ である。
次数2以下の純被約連結曲線の種数は零なので、双方とも次数3に限られる。
定理\ref{thm:pure-curve-defect}を使う。
すなわち次数6なら全純部分の種数3に対し二つの開閉楕円種数成分の和の算術種数は1、
次数7なら残る純直線を加えた種数は0で所要の5に達しない。
従って種数一の連結成分は保たれる（同じ議論は節点の輪にも適用できる）。

この唯一の特異点の環を明記する。$A\subset A'$ の余長が1なら $\mathfrak m_A A'\subset A$、
定数項が零なので $\mathfrak m_A$ は $A'$ のイデアルである。
$\mathfrak m_{A'}/\mathfrak m_A$ は一次元の加群なので $\mathfrak m_{A'}^2\subset\mathfrak m_A$。
従って定数を貼ったうえに一階係数に一つの線型条件を課すだけである。
係数の零となる枝は残りと値だけを貼る滑らかな軸である。この軸上の余接捩れは残りの環の接次元に等しく、高々2である。
実際、残りの最小座標 $y_j$ と軸の $t$ に対し追加の関係は $t y_j=0$ であり、残りの関係には一次項がない。
よってこの枝で恒等なら必要な長さ4を得られず、非恒等でも $p$ の固定により少なくとも3の捩れが必要である。

全係数が非零なら各枝の変数の倍率を変えて一階係数の和を零とできる。
$r=1,2$ は尖点と二枝接点であり、形式表示はそれぞれ $y^2=x^3$、$y(y-x^2)=0$ である。
$r\ge3$ ならこの環は $r-1$ 次元空間中の $r$ 本の直線の錐で、方向は射影枠をなす。
実際、和零超平面の成分ごとの積は各座標の二次項を張る
（相異なる $i,j,k$ に対し $(e_i-e_j)(e_i-e_k)=e_i$）。
以後全ての高次項を生成でき、接次元から $r\le4$ である。
$r=3$ は平面三重点、$r=4$ は平面内の枠四点を切る二つの二次式の錐である。
尖点では $(\deg B,\deg T)=(1,3)$、他の三型の各枝では $(0,2)$ である。
接点と三重点では平面方程式の勾配の共通位数を取る。四枝では各方向で二つの勾配は独立で、
枝でともに一位となる。いずれも $\deg M+\deg T=0$、$\deg M<0$。
四移動が枝上で恒等なら $2\deg M+\deg T=0$ に反し、それ以外では一致長零に固定点 $p$ が反する。
\end{proof}

\begin{theorem}\label{thm:pentagon-node-rank}
点ファイバーを持つ標準有限Frobenius窓の点グラフは、開閉な空間五角形を含まない。
特に被約な種数一の純部分を持つ五次点グラフの標準AS窓は滑らかな楕円五次の二類だけである。
\end{theorem}
\begin{proof}
五頂点を隣接順に番号付けし、全位置でグラフ対応により同じ番号を使う。
各位置で任意の四頂点は独立である。四点は巡回的な連続四点であり、
同一平面に入ればその両端の非隣接二辺が交わってしまう。
従って各位置で頂点を同じ射影枠 $v_0,\ldots,v_4$ に取り、代表を $\sum v_j=0$ とできる。
各直線上の隣接埋込み間の写像は
$[s v_j+t v_{j+1}]\mapsto[s v_j+q_j t v_{j+1}]$、$q_j\ne0$ である。
五倍率の一致や添字方向の周期を仮定しない。

この辺の関係 $B$ に対し $N_{jk}=B(v_j,v_k)$ は行和・列和が零で、
$N_{jj}=0,\ N_{j+1,j}=-q_j N_{j,j+1}$ を満たす。
逆にこれらの行列が六次元をなす。隣接五成分 $N_{j,j+1}$ を任意に取ると、
非隣接の十成分は行と列を頂点とする二部十角形の辺に対応する
（行 $j$ は列 $j\pm2$ に接する）。各頂点での和を指定する条件は一つだけ従属し、
行和総計と列和総計の一致で可解である。従って六関係はこの空間全体である。

$U_i$ をこの空間で行 $i$ が零の部分とする。片側点ファイバー条件から次元3である。
$U_0$ の列 $j$ での評価階数は $j=0$ で高々1、他の四点でちょうど2である。
計算を示す。$a,b,c=N_{12},N_{23},N_{34}$ を自由変数とする。行0が零なら
残りの非隣接成分は和条件から一意である（十角形を行0で切った木で和を順に解く）。
各列のうち行 $1,2,3$ を取れば十分で、その $(a,b,c)$ からの係数行列は順に
\[
\begin{gathered}
\begin{pmatrix}0&0&0\\q_1-1&0&1-q_3\\1-q_1&0&q_3-1\end{pmatrix},\quad
\begin{pmatrix}0&0&0\\-q_1&0&0\\q_1-1&q_2&-q_3\end{pmatrix},\quad
\begin{pmatrix}1&0&0\\0&0&0\\0&-q_2&0\end{pmatrix},\\
\begin{pmatrix}0&-1&q_3\\0&1&0\\0&0&0\end{pmatrix},\quad
\begin{pmatrix}-1&1&-q_3\\1&-1&q_3-1\\0&0&1\end{pmatrix}.
\end{gathered}
\]
最後の四行列には順に $-q_1q_2,-q_2,-q_3,1$ の二次小行列がある。
転置した辺では倍率を逆数にすれば同じ結論となる。

四次テンソル $w$ の位置0で節点 $p$ を評価する。位置1と2の関係を用いると
残る三次テンソルは $U_p\otimes V_3$ に入り、$V_3^*$ の縮約像は $U_p$ 全体である。
実際 $R_3$ と $R_1$ の完全対合のもとでこれは位置0での $R_3$ の評価の転置に当たり階数3である。
同様に位置2と3の関係を転置して使えば、位置3の前点 $j=\sigma^{-1}p$ での列が零になる
三次元部分全部を与える（位置3の消滅超平面は最後から最初への辺で決まる）。
従って位置2でも節点 $z$ を評価した $w(p,-,z,-)$ の階数は、
第一の表示では $z=p$ のときだけ1以下、第二では $z=j$ のときだけ1以下である。
よって全節点で $\sigma^{-1}p=p$ を要し、前定理の固定点を持たない回転に反する。
次数上限7から五角形自体は四移動で保たれる。

最後に五次で純部分が被約な種数一なら定理\ref{thm:pure-curve-defect}より有限部分は零である。
非連結なら次数2以下の被約な連結成分があり、種数零で除外される。
多角形の回転の位数 $k>1$ は成分の全軌道の長さであり、各軌道の成分次数は等しい。
よって $k\mid 5$、成分は五本の直線に限られ、今排除した。整の場合は
定理\ref{thm:elliptic-quintic-regularity}の二類で尽きる。
\end{proof}
